\begin{claim}{C.39}{Proved}{A global contraction for loops with commuting values}
Let $f_t$, $0\leq t\leq1$, be a based loop in
$\operatorname{Homeo}^+(\R^2)$ with
$f_s f_t=f_t f_s$ for every $s,t$.
Suppose the loop is null-homotopic in that ambient group.
For every subgroup $H$ commuting with all $f_t$, the loop
has a based null-homotopy in
$C_{\operatorname{Homeo}^+(\R^2)}(H)$.
The homotopy and its inverse family are jointly continuous
on the whole plane, including every change and accumulation
of orbit types.

Consequently, on every connected admissible surface $U$
of C.1, an element of the kernel of the exact inclusion map
\[
 \pi_1 C_{\operatorname{Homeo}(U)}(H)
      \longrightarrow\pi_1\operatorname{Homeo}(U)
                                                        \tag{C.39.1}
\]
that is represented by a loop with pairwise commuting values
is zero. On a hyperbolic surface every such based loop in
$C_{\operatorname{Homeo}(U)^0}(H)$ is therefore null-homotopic
in that centralizer.
\end{claim}

\begin{proof}
\emph{The actual image on every orbit.}
Set $G=\langle f_t\rangle$. This is an abelian group.
The orbit classification C.28 applies.
On an open orbit $O$, the closure $K_O$ of $G|_O$ in
$\operatorname{Homeo}(O)$ is Polish, abelian, connected,
and transitive. It acts freely: an element fixing a point
commutes with the transitive subgroup $G|_O$ and hence fixes
every point. The Polish open orbit-map theorem identifies
evaluation $K_O\to O$ homeomorphically, with trivial
stabilizer. Thus $K_O$ is locally compact. Its faithful
action on $O\times\R$ and the Hilbert--Smith input of C.33
make it a two-dimensional abelian Lie group.
Since $O$ is a noncompact open subset of the plane, the
possibilities are $\R^2$ and $\R\times S^1$.
The original image is already the whole group:
for each $k\in K_O$, transitivity of $G|_O$ gives an
element with the same evaluation, and freeness makes
the two elements equal.

On an intrinsic line or circle orbit $B$, C.28 gives
compact-open continuity of the restricted original paths.
The same abelian closure and free evaluation argument,
now in the intrinsic topology, identifies $G|_B$ with
the full translation group $\R$ or rotation group $S^1$.
This is also the abelian case of (C.37.9).
On a point orbit the image is trivial.
Denote the resulting actual Lie image on any orbit $B$
by $K_B$. Its simply connected covering group is a vector
group of dimension at most two.

The original restricted loop has a closed based lift
$\widetilde n_B(t)$ to this vector group.
For an open orbit, C.32 gives its null-homotopy in
$\operatorname{Homeo}(O)$; evaluation and the
homeomorphism $K_O\cong O$ then make its class in
$\pi_1K_O$ zero. For a line there is no covering kernel.
For a circle the based lift closes by the exact winding
comparison C.29. Define on each orbit
\[
              R^B_{r,t}
                =p_B\bigl(r\,\widetilde n_B(t)\bigr),
                    \qquad 0\leq r,t\leq1.
                                                        \tag{C.39.2}
\]
Here $p_B$ is the Lie-group covering homomorphism,
with the trivial interpretation for a point orbit.
The action of the right side on $B$ is its actual
homeomorphism action. Its endpoints are the identity
at $r=0$ and the original loop at $r=1$; both based
edges $t=0,1$ are fixed.

\emph{One compact word family for all orbit types.}
In dimension two, C.37 represents every value of
(C.39.2) by a word isotopy
\[
 W_{\epsilon,a,b,t}(u)
        =f_{ut}^{\,\epsilon}f_{ub}f_{ua}^{-1},
 \quad \epsilon\in\{0,1\},\quad 0\leq a,b\leq t\leq1,
                                                        \tag{C.39.3}
\]
with equality of its specified based universal-cover
endpoint. In dimension one the intermediate value
theorem already finds $b\leq t$ with
$\widetilde n_B(b)=r\widetilde n_B(t)$.
This is a member of (C.39.3) with $\epsilon=0,a=0$.
The identity word handles point orbits.
Thus the same compact family of original plane isotopies
works for every orbit, with the based lifts retained.

Take arbitrary points $x_i\to x$ in the plane and
parameters $(r_i,t_i)\to(r,t)$, and write $B_i=Gx_i$,
$B=Gx$. Choose representing word isotopies $W_i$ for
$R^{B_i}_{r_i,t_i}$. Compactness of the parameter set
in (C.39.3) gives a subsequence converging, together
with inverses and uniformly in $u$, to a word isotopy
$W$. Put $w=W(1)$.
If $B$ is open, all sufficiently large $x_i$ are in
that same orbit and continuity follows from its
Lie action. If $B$ is a point, every possible word
limit fixes $x$, giving the required convergence.
It remains to identify the limiting word on a line
or circle. This argument explicitly allows the $B_i$
to have different types.

For each fixed positive integer $n$, set
$m_i=\lfloor nr_i\rfloor$. In the vector covering
group of the actual image $K_{B_i}$, one has
\[
 n\,\widetilde w_i
       -m_i\,\widetilde n_{B_i}(t_i)
       =(nr_i-m_i)\widetilde n_{B_i}(t_i).
                                                        \tag{C.39.4}
\]
The right side is represented by another word isotopy
$Z_i$ from (C.39.3), with equality of based lifted
endpoints. For this fixed $n$, pass to a further
subsequence so that $m_i=m_n$ is constant and
$Z_i\to Z_n$. Then $|m_n-nr|\leq1$.
Identities on $B_i$ pass to the whole orbit $B$ by
evaluation at $hx_i\to hx$, for any fixed $h\in G$.
In particular, with $g=f_t$ and $z_n=Z_n(1)$,
\[
             (w|_B)^n(g|_B)^{-m_n}=z_n|_B.
                                                        \tag{C.39.5}
\]

For a line, let $\theta(u)$ be the real translation
parameter of $f_u|_B$, and put
$M=\max_u|\theta(u)|$.
The real parameters $A$ of $w|_B$ and $C_n$ of
$z_n|_B$ come from their word paths:
each has the form
$\epsilon\theta(t)+\theta(b)-\theta(a)$.
Consequently $|C_n|\leq3M$, and (C.39.5) is the
real identity $nA-m_n\theta(t)=C_n$.
Division by $n$ gives $A=r\theta(t)$, as required.

For a circle, use the based real rotation lift
$\theta$ of the original loop, so $\theta(0)=\theta(1)=0$.
The word paths again specify real endpoint angles
$A,C_n$ with $|C_n|\leq3M$, but (C.39.5) initially
gives their relation only modulo integers.
The exact lift in (C.39.4) removes this ambiguity.
Indeed,
\[
       \Lambda_i(u)
         =W_i(u)^n f_{ut_i}^{-m_i}Z_i(u)^{-1}
                                                        \tag{C.39.6}
\]
is a based loop in the actual group $K_{B_i}$ and
has a closed based universal-cover lift.
It is therefore null-homotopic in that group.
Choose $h\in G$ with $hx\ne x$.
For large $i$ the two points $x_i,hx_i$ of $B_i$
are distinct. Evaluation of (C.39.6) at them has
difference-direction winding zero in
$\operatorname{Conf}_2(\R^2)$, regardless of whether
$B_i$ is an open, line, or circle orbit.
This follows from the actual null-homotopy in
$K_{B_i}$ and its continuous evaluation into the
plane; C.28 supplies that evaluation for intrinsic
line orbits.

The word isotopies and the two evaluation points
converge. Their limiting pair stays distinct.
The winding therefore remains zero. On the limiting
Jordan circle its value, by C.29, is
$nA-m_n\theta(t)-C_n$, up to a common orientation
sign. Thus the real identity
\[
 nA-m_n\theta(t)=C_n,\qquad
 |A-r\theta(t)|
       \leq\frac{3M+|\theta(t)|}{n}
                                                        \tag{C.39.7}
\]
holds for every $n$. It forces $A=r\theta(t)$.
This also treats an identity endpoint rotation
whose based angle $\theta(t)$ is a nonzero integer.
All possible subsequential word limits therefore
give the prescribed value (C.39.2).

Compactness of (C.39.3) now proves joint continuity
of (C.39.2) on the entire plane.
The inverse loop has the same orbits and image
groups, and its construction gives precisely the
inverse of (C.39.2); the same proof gives joint
inverse continuity. Each map is an orbitwise
bijection, so these are plane homeomorphisms.
Every commuting operator carries the actual
abelian Lie image on one orbit isomorphically to
that on the transported orbit. These are continuous
Lie-group isomorphisms in the specified intrinsic
or open-orbit topology. Their based covering lifts
are linear, and hence commute with multiplication
by $r$. Formula (C.39.2) is therefore exactly
$H$-equivariant. Taking $r=1-s$ proves the plane
assertion.

\emph{The precise surface lift and descent.}
Every connected admissible surface in C.1 has
universal covering $p:\widetilde U\to U$ with
$\widetilde U\cong\R^2$. Let $f_t$ represent an
element of the kernel in (C.39.1), with its
commuting-value property. Choose an ambient
filling $A:D^2\to\operatorname{Homeo}(U)$ whose
boundary is this loop and whose specified based
boundary point maps to the identity.
The actual map
\[
 D^2\times\widetilde U\longrightarrow U,\qquad
                  (z,x)\longmapsto A(z)(p(x))
                                                        \tag{C.39.8}
\]
has a unique lift equal to the identity on the
specified based fiber. Its domain is simply
connected. Lifting the inverse family as well,
and using uniqueness, proves that its slices are
homeomorphisms with jointly continuous inverses.
It is an ambient filling of a closed based
plane loop $\widetilde f_t$.

For every lift $\widetilde h$ of an operator,
$[\widetilde f_t,\widetilde h]$ is a deck
transformation. It depends continuously on $t$
and starts at the identity, so it is identically
the identity. Similarly
$[\widetilde f_s,\widetilde f_t]$ is a deck
transformation depending continuously on $(s,t)$;
it is the identity on the axes, and hence
everywhere. Thus the lifted loop has commuting
values and centralizes the full group of
operator lifts, including the deck group.
Apply the plane assertion. Its equivariant
contraction descends through $p$, with both
joint continuities, to the original centralizer
on $U$. Its slices lie in the identity component
by the displayed paths. This proves (C.39.1).
For a hyperbolic $U$ the ambient identity
component is contractible, so the ambient
kernel condition is automatic.
\end{proof}

\paragraph{Scope of the global subcase.}
This theorem treats loops with commuting values on
every orbit at once, including open orbits whose
operator image has dimension zero or one.
It does not identify every fundamental-group class
of an arbitrary centralizer with such a loop.
No deformation of a general centralizer loop to
this subcase has been constructed here.
C.12 remains open for general loops.

\paragraph{The earlier abelian closed subgroup cannot be a centralizer.}
Fix the specific embedding $j:B\hookrightarrow A^0$ of C.26, where
$B=L^1([0,1],\Z)/\Z$, $A=\operatorname{Homeo}(U)$ and $U$ is
hyperbolic. For every $H\leq A$ commuting with $j(B)$, the induced map
$\pi_1(B)\to\pi_1(C_A(H),1)$ is zero. Indeed the displayed generator
of $\pi_1(B)=\Z$ has pairwise commuting values and is ambient-null,
so C.39 contracts its image; the induced homomorphism therefore
vanishes on the whole group. The based contraction lies in $A^0$,
so the same conclusion holds with target $\pi_1(C_{A^0}(H),1)$.
Consequently $j(B)$ cannot equal either of these centralizers,
since its topological-group identification with $B$ induces an
isomorphism on the nonzero fundamental group. This rules out that
exact realization as a zero-level counterexample.
